\section{Common Training Contract}
\label{sec:common}

\subsection{Scope and evidence}

All five ideas begin with an offline personalized language-model policy and
continue learning from chronological agent interactions.  The common policy
state is
\begin{equation}
  \params = \bigl(\phi,e_0,\{e_u:u\in\users\}\bigr),
  \label{eq:common-state}
\end{equation}
where the backbone is frozen, $\phi$ is a shared LoRA adapter, $e_0$ is a
generic soft prompt, and $e_u$ is a private embedding for user $u$.
This organization follows personalized RLHF and parameter-efficient
fine-tuning \citep{li2024personalized,hu2022lora}.  Preference losses use DPO
when a pairwise label is available \citep{rafailov2023dpo}; numerical updates
use AdamW \citep{loshchilov2019adamw}.

\Evidence.  The cited papers establish prior algorithms and motivate the
proposed transformations.  They do not establish that any method in this note
works.  Every improvement claim therefore remains a hypothesis until a
checkpoint-changing run and untouched future evaluation have been completed.

\subsection{Shared record and checkpoint schema}

Every trainable record must identify the user and profile version, task and
environment, generating model, behavior checkpoint, rendered state, actions
and tool calls, outcome or preference provenance, timestamp, data split, and
training eligibility.  Token log-probabilities and replayable state
fingerprints are required when an algorithm uses them.  Missing fields cause
quarantine; they are never silently reconstructed.

Each candidate checkpoint belongs to an append-only directed acyclic graph.
Its receipt contains the parent hash, adapter hash, optimizer-state hash,
data-policy version, evaluation receipt, activation status, and rollback
reason.  A rejected candidate cannot generate later data or donate optimizer
moments.

\subsection{Shared lifecycle}

\begin{table}[t]
  \centering
  \caption{Lifecycle used by all five workflows.  The middle step changes by
  idea; the provenance and activation steps do not.}
  \label{tab:shared-lifecycle}
  \begin{tabularx}{\textwidth}{@{}L{0.09\textwidth}Y Y@{}}
    \toprule
    Phase & Required action & Output \\
    \midrule
    Offline & Train the registered \prlhf\ or \agenticdpo\ parent on a
    chronological historical split. & Immutable checkpoint $c_0$, fixed
    reference policy, optimizer receipt. \\
    Collect & Run only the active checkpoint in the declared environment and
    validate every new record. & Versioned interaction ledger. \\
    Adapt & Execute the idea-specific update without reading validation or
    future-test evidence. & Candidate $c_{t+1}$ and complete update receipt. \\
    Verify & Load $c_{t+1}$ in a fresh worker and compare it with $c_t$ on a
    fixed validation panel. & Activation or rollback decision. \\
    Evaluate & Test activated checkpoints on untouched future periods from
    natural initial states. & Separate task-success and personalization
    estimates. \\
    \bottomrule
  \end{tabularx}
\end{table}

\subsection{Upper, middle, and lower algorithm layers}

The \emph{upper layer} names the optimized object: the shared LoRA and
user-conditioned parameters in \cref{eq:common-state}.  The \emph{middle
layer} maps evidence to an update, such as preference optimization, replay,
policy gradient, bilevel weighting, or active acquisition.  The \emph{lower
layer} provides loops, buffers, projections, checkpointing, and stopping
rules.  A method is credited only for the layer it changes.  Standard
minibatching, AdamW, replay storage, and checkpoint hashing remain
infrastructure.

\subsection{Common activation gate}
\label{sec:activation-gate}

A candidate may become active only if all of the following hold:
\begin{enumerate}
  \item its adapter and optimizer hashes load in a clean worker;
  \item the update is finite and nonzero;
  \item task success and safety are non-inferior on the fixed validation panel;
  \item the declared personalization validation statistic improves;
  \item no protected provenance or future-test boundary was crossed.
\end{enumerate}
Failure restores both parent parameters and parent optimizer state.  Runs stop
when their fixed interaction, token, optimizer-step, GPU, or wall-clock budget
is exhausted; after three consecutive windows without an activated candidate;
or at a preregistered safety termination.
